%==============================================================================
\section{Phân cụm dựa trên phân phối (Distribution-Based)}
\label{sec:distribution-based}
%==============================================================================

\begin{dinhnghia}[title={Phân cụm dựa trên phân phối}]
Còn gọi \textbf{phân cụm xác suất}: giả định dữ liệu sinh từ \textbf{hỗn hợp} phân phối xác suất;
ước lượng phân phối nền rồi gom điểm có tính chất tương tự.
\end{dinhnghia}

\subsection{Gaussian Mixture Model (GMM)}

\begin{dinhnghia}[title={GMM}]
Giả định dữ liệu từ hỗn hợp phân phối \textbf{Gaussian}; mỗi Gaussian là một cụm.
Ưu điểm: cụm \textbf{chồng lấn}, mô hình hóa \textbf{hiệp phương sai}, gán cụm \textbf{xác suất} cho từng điểm.
Ứng dụng: phân đoạn ảnh, nhận dạng mẫu, phát hiện bất thường.
\end{dinhnghia}

\subsection{Triển khai Python (GaussianMixture, covariance full)}

\begin{vidu}[title={Ví dụ GMM}]
\begin{lstlisting}
from sklearn.mixture import GaussianMixture
from sklearn.datasets import make_blobs
import matplotlib.pyplot as plt

X, _ = make_blobs(n_samples=200, centers=4, random_state=0)

gmm = GaussianMixture(n_components=4, covariance_type='full')
gmm.fit(X)
labels = gmm.predict(X)

print("Cluster labels:", labels)
plt.figure(figsize=(7.5, 3.5))
plt.scatter(X[:, 0], X[:, 1], c=labels, cmap='viridis')
plt.show()
\end{lstlisting}
\end{vidu}

\begin{ynghia}[title={covariance\_type}]
\texttt{full}: ma trận hiệp phương sai đầy đủ mỗi cụm.
Các lựa chọn khác: \texttt{tied}, \texttt{diag}, \texttt{spherical}.
Tham số khởi tạo điều khiển cách ước lượng tham số Gaussian ban đầu.
\end{ynghia}

\tsimg{13_distribution_based/img01.jpg}

\subsection{Ưu và nhược điểm GMM}

\begin{tinhchat}[title={Ưu điểm}]
\begin{itemize}
  \item Mô hình hóa phân phối \textbf{linh hoạt}.
  \item Xử lý dữ liệu \textbf{thiếu}.
  \item Khung \textbf{xác suất} (độ không chắc chắn gán cụm).
  \item Ước lượng mật độ và \textbf{sinh mẫu} mới.
  \item Dùng trong \textbf{bán giám sát} khi một phần có nhãn.
\end{itemize}
\end{tinhchat}

\begin{luuy}[title={Nhược điểm}]
\begin{itemize}
  \item \textbf{Nhạy} số cụm và khởi tạo mean/covariance.
  \item \textbf{Đắt} trên chiều cao (nghịch đảo ma trận hiệp phương sai).
  \item Giả định \textbf{Gaussian} có thể không đúng với mọi dữ liệu.
  \item Dễ \textbf{overfitting} khi nhiều tham số hoặc ít mẫu.
  \item Khó diễn giải khi ma trận hiệp phương sai phức tạp.
\end{itemize}
\end{luuy}
